\documentclass[10pt,a4paper]{amsart}
\usepackage[margin=19mm]{geometry}
\usepackage{amsmath,amssymb,mathtools}
\usepackage[scr=rsfs,cal=euler]{mathalfa}
\usepackage{xfrac}
\usepackage[unicode,hidelinks,bookmarks=false]{hyperref}
\newcommand{\ie}{i.e.\ }
\newcommand{\cf}{cf.\ }
\newcommand{\resp}{respectively\ }
\newcommand{\Subsection}{\subsection}
\providecommand{\nobreakdash}{\nobreak\hbox{-}}
\setlength{\emergencystretch}{2em}
\multlinegap0pt
\usepackage{bm}
\usepackage{bbm}
\usepackage{dsfont}
\usepackage{xspace}
\usepackage{cancel}
\usepackage{mathtools}
\usepackage{amsthm}
\makeatletter
\@ifundefined{thm}{\newtheorem{thm}{Theorem}}{}
\makeatother
\def\vol{{\rm vol}}
\newcommand{\R}{\mathbb R}
\newcommand{\Z}{\mathbb Z}
\newcommand{\cV}{\mathcal V}
\title[On the Log-submodularity for zonoids: from Mixed Volume ineq]{On the Log-submodularity for zonoids: from Mixed Volume inequalities to the Hypercube}
\author{Gennadiy Averkov, Katherina von Dichter, Ivan Soprunov}
\date{Source: arXiv:2608.14909v1 (CC BY 4.0)}
\begin{document}
\maketitle

\noindent\emph{Source and licence.} On the Log-submodularity for zonoids: from Mixed Volume inequalities to the Hypercube. Version arXiv:2608.14909v1 (CC BY 4.0), distributed under the Creative Commons Attribution 4.0 International licence, as recorded in the arXiv licence metadata for this exact version and shown on the abstract page https://arxiv.org/abs/2608.14909v1. Full rights notice: licenses/arxiv-2608.14909-CC-BY-4.0.txt.\par\smallskip

\noindent\textit{Editorial scope.} Counted unit: the Thm whose source label is \texttt{None} in the article (source0.tex lines 1022--1025, source label \texttt{None}), delivered together with the article's own complete original proof of it (source0.tex lines 1027--1052, one \texttt{proof} environment of 26 lines). No mathematics is written, repaired or re-derived: statement and proof are the article's own text line for line. Required context retained uncounted: e:vol-poly (lines 891--893); e:Grass-dual (lines 1017--1019). Every definition, display and label of those lines is retained unchanged. Omitted from this package and not redistributed: 1--890, 894--1016, 1020--1021, 1026--1026, 1053--2085. Nothing inside a retained excerpt is omitted or altered: every retained body is delivered exactly as the article's lines are written, so no omission receipt of any kind accompanies this package. Citations: the retained excerpts carry no citation command at all, so no bibliography entry of the article is transcribed and no reference is redistributed. Reference adaptation: none was needed and none was made; every reference inside the delivered unit resolves inside it, so no unresolved reference or citation is produced. Printed slips kept literal: no printed word, symbol, hypothesis or number of the counted statement or of its proof was altered. Layout and class: the article's own class is \texttt{amsart} with packages including graphicx, tikz, latexsym, color, amsmath, systeme, amsthm, amssymb, amscd, amsfonts, todonotes, hyperref, pgf, pgfplots; the wrapper is the lane's pinned \texttt{amsart} editorial wrapper at 10pt on A4, which restates the theorem-like environments the retained text needs and adds the article's own macro definitions that the retained text needs (\texttt{\textbackslash R}, \texttt{\textbackslash Z}, \texttt{\textbackslash cV}, \texttt{\textbackslash vol}), with extra wrapper packages (bm, bbm, dsfont, xspace, cancel, mathtools). The delivered numbering is the unit's own. Assets: no retained span contains a figure environment, a graphics-inclusion command, a PDF-page inclusion, a table or a TikZ picture; no image file is copied into this package, the package declares no asset dependency and no image file is redistributed with it. Licence: version arXiv:2608.14909v1 of this article is distributed under the Creative Commons Attribution 4.0 International licence, as recorded in the arXiv licence metadata for this exact version and shown on the abstract page https://arxiv.org/abs/2608.14909v1. A full rights notice is delivered at licenses/arxiv-2608.14909-CC-BY-4.0.txt. Characters: every retained excerpt is pure ASCII, so the delivered engine, XeLaTeX with its default Latin Modern text font, reports no missing character.
\begin{align}\label{e:vol-poly}
 f_{A,B}(x)= \sum_{p=0}^k\sum_{|J|=p}\frac{n!}{(n-p)!}V(A[n-p],B_J)\,x_{J}=\sum_{p=0}^k\sum_{|J|=p}\vol(P_{[B_{J}]^\perp}A)\,x_{J},
\end{align}

\begin{equation}\label{e:Grass-dual}
p_{I^c}(L^\perp)=\pm p_I(L),   
\end{equation}

\begin{thm}
The map $f(x_1,\dots, x_k)\mapsto x_1\cdots x_k f(x_1^{-1},\dots, x_k^{-1})$ defines a bijection between the
mixed volume diagrams $\cV(n,m,k)$ and $\cV(m+k-n,m,k)$.
\end{thm}

\begin{proof} 
Let $f_{A,B}(x)=\vol_d(A+x_1B_1+\dots+x_kB_k)$ be the volume polynomial of $A\in\Z_m^n$ and $B_1,\dots, B_k\in\Z_1^n$. Without loss of generality, we may assume that $B_1,\dots, B_k$ are pairwise orthogonal unit segments.

Form a matrix $M\in\R^{n\times (m+k)}$ whose columns are the generators of $A$ and $B_1,\dots, B_k$ and let $L$ be the row span of $M$. Let $M^\perp\in\R^{(m+k-n)\times(m+k)}$ be a matrix whose row span is the orthogonal complement $L^\perp$ and let $A^*\in\Z_m^{m+k-n}$ be the zonotope generated by the first $m$ columns of $M^\perp$ and $B^*_1,\dots, B^*_k\in\Z_1^{m+k-n}$ be segments generated by the last $k$ columns of $M^\perp$. Recall that by (\ref{e:vol-poly})
$$
f_{A,B}(x)=\sum_{p=0}^k\sum_{|J|=p}\vol(P_{[B_{J}]^\perp}A)\,x_{J}.
$$
Note that the volume of the projection $P_{[B_{J}]^\perp}A$ equals the sum of the absolute values of the maximal minors that use exactly those of the last $k$ columns of $M$ that correspond to the subset $J\subset[k]$. Thus, we can write
$$
\vol(P_{[B_{J}]^\perp}A) = \sum_{S\subset I\subset T}|p_I(L)|,
$$ 
where the sum is over subsets $I\subset[m+k]$ of size $n$, $S=\{m+i : i\in J\}$, and $T=[m]\cup S$.

Applying the duality (\ref{e:Grass-dual}), we get (up to some positive scalar)
$$
\sum_{S\subset I\subset T}|p_I(L)| = \sum_{T^c\subset I^c\subset S^c}|p_{I^c}(L^\perp)|,
$$
where the complements are taken with respect to the set $[m+k]$. Note that $T^c=\{m+i : i\in [k]\setminus J\}$ and $S^c=[m]\cup T^c$.
Therefore, 
$$
\sum_{T^c\subset I^c\subset S^c}|p_{I^c}(L^\perp)|=\vol(P_{[B^*_{[k]\setminus J}]^\perp}A^*).
$$
Comparing this with the formula for the volume polynomial (\ref{e:vol-poly}) and noting $x_J=x_1\cdots x_k(x_{[k]\setminus J})^{-1}$ we obtain (up to a positive scalar)
$$f_{A,B}(x_1,\dots, x_k)=x_1\cdots x_k f_{A^*,B^*}(x_1^{-1},\dots, x_k^{-1}),$$
which completes the proof.
\end{proof}

\end{document}
